\textbf{Phần 1. Trắc nghiệm lựa chọn}
\setcounterpageref{ex}{0}

\begin{ex} %[3D16N81]
Số liền sau của số $89\,999$ là
\choice
{\True $90000$}
{$89998$}
{$90001$}
{$100000$}
\loigiai{
$89999 + 1 = 90000$.\\
\textbf{Đáp án của câu hỏi là A}
}
\end{ex}

\begin{ex} %[3D16N81]
Chữ số La Mã $\text{XX}$ biểu thị số nào?
\choice
{$10$}
{\True $20$}
{$15$}
{$25$}
\loigiai{
$\text{XX} = 20$.\\
\textbf{Đáp án của câu hỏi là B}
}
\end{ex}

\begin{ex} %[3D16H81]
Kết quả của phép tính $45\,000 + 35\,000$ là
\choice
{$70000$}
{$90000$}
{$75000$}
{\True $80000$}
\loigiai{
$45000 + 35000 = 80000$.\\
\textbf{Đáp án của câu hỏi là D}
}
\end{ex}

\begin{ex} %[3D16H81]
Kết quả của phép tính $100\,000 - 45\,000$ là
\choice
{$45000$}
{\True $55000$}
{$65000$}
{$50000$}
\loigiai{
$100000 - 45000 = 55000$.\\
\textbf{Đáp án của câu hỏi là B}
}
\end{ex}

\begin{ex} %[3D16H81]
Kết quả của phép nhân $15\,000 \times 4$ là
\choice
{$50000$}
{$70000$}
{\True $60000$}
{$80000$}
\loigiai{
$15000 \times 4 = 60000$.\\
\textbf{Đáp án của câu hỏi là C}
}
\end{ex}

\begin{ex} %[3D16H81]
Kết quả của phép chia $84\,000 : 4$ là
\choice
{$20000$}
{$22000$}
{$24000$}
{\True $21000$}
\loigiai{
$84000 : 4 = 21000$.\\
\textbf{Đáp án của câu hỏi là D}
}
\end{ex}

\begin{ex} %[3H16V81]
Một hình vuông có chu vi $48\text{ cm}$. Diện tích hình vuông đó là
\choice
{$96\text{ cm}^2$}
{$120\text{ cm}^2$}
{\True $144\text{ cm}^2$}
{$160\text{ cm}^2$}
\loigiai{
Cạnh: $48 : 4 = 12\text{ cm}$.\\
Diện tích: $12 \times 12 = 144\text{ cm}^2$.\\
\textbf{Đáp án của câu hỏi là C}
}
\end{ex}

\begin{ex} %[3H16V81]
Một hình chữ nhật có chiều dài $15\text{ cm}$, chiều rộng $8\text{ cm}$. Diện tích hình chữ nhật là
\choice
{$110\text{ cm}^2$}
{\True $120\text{ cm}^2$}
{$130\text{ cm}^2$}
{$100\text{ cm}^2$}
\loigiai{
$15 \times 8 = 120\text{ cm}^2$.\\
\textbf{Đáp án của câu hỏi là B}
}
\end{ex}

\begin{ex} %[3D16V81]
Mua 3 quyển sách giá $25\,000$ đồng/quyển và đưa tờ $100\,000$ đồng. Người bán hàng phải trả lại
\choice
{\True $25000$ đồng}
{$20000$ đồng}
{$30000$ đồng}
{$35000$ đồng}
\loigiai{
$100000 - 25000 \times 3 = 25000$ đồng.\\
\textbf{Đáp án của câu hỏi là A}
}
\end{ex}

\begin{ex} %[3D16V81]
Làm tròn số $67\,850$ đến hàng nghìn ta được
\choice
{$67000$}
{$70000$}
{\True $68000$}
{$67900$}
\loigiai{
Hàng trăm là 8 nên làm tròn lên $68000$.\\
\textbf{Đáp án của câu hỏi là C}
}
\end{ex}

\begin{ex} %[3H16V81]
Tháng 5 có bao nhiêu ngày?
\choice
{\True $31$ ngày}
{$30$ ngày}
{$28$ ngày}
{$29$ ngày}
\loigiai{
Tháng 5 có 31 ngày.\\
\textbf{Đáp án của câu hỏi là A}
}
\end{ex}

\begin{ex} %[3D16C81]
Một khu đất hình chữ nhật có chu vi $100\,000\text{ m}$, chiều dài gấp 4 lần chiều rộng. Chiều dài khu đất đó là
\choice
{$35000\text{ m}$}
{$45000\text{ m}$}
{$50000\text{ m}$}
{\True $40000\text{ m}$}
\loigiai{
Nửa chu vi: $100000 : 2 = 50000\text{ m}$.\\
Chiều dài: $50000 : 5 \times 4 = 40000\text{ m}$.\\
\textbf{Đáp án của câu hỏi là D}
}
\end{ex}

\textbf{Phần 2. Trắc nghiệm đúng sai}
\setcounterpageref{ex}{0}

\begin{ex} %[3D16H81]
Xét tính đúng sai của các kiến thức tổng hợp:
\choiceTF
{\True $1\text{ m} = 1000\text{ mm}$.}
{\True Diện tích hình vuông cạnh $10\text{ cm}$ là $100\text{ cm}^2$.}
{Tháng 4 có 31 ngày.}
{\True Tháng 4 có đúng 30 ngày.}
\loigiai{
\begin{itemize}
    \item $1\text{ m} = 1000\text{ mm}$: Đúng.
    \item $10 \times 10 = 100\text{ cm}^2$: Đúng.
    \item Tháng 4 có 30 ngày (không phải 31). Ý 3 Sai, Ý 4 Đúng.
\end{itemize}
\begin{itemize}
    \item a) \textbf{Đúng}.
    \item b) \textbf{Đúng}.
    \item c) \textbf{Sai}.
    \item d) \textbf{Đúng}.
\end{itemize}
}
\end{ex}

\begin{ex} %[3D16V81]
Một cửa hàng có $90\,000\text{ kg}$ xi măng. Lần 1 bán được $35\,000\text{ kg}$, lần 2 bán được $40\,000\text{ kg}$. Xét tính đúng sai:
\choiceTF
{\True Cả 2 lần cửa hàng bán được $75\,000\text{ kg}$ xi măng.}
{\True Lượng xi măng còn lại là $15\,000\text{ kg}$.}
{Lần 2 bán ít hơn lần 1 là $5\,000\text{ kg}$.}
{\True Lần 2 bán nhiều hơn lần 1 là $5\,000\text{ kg}$ ($40000 - 35000 = 5000$).}
\loigiai{
\begin{itemize}
    \item Đã bán: $35000 + 40000 = 75000\text{ kg}$. Đúng.
    \item Còn lại: $90000 - 75000 = 15000\text{ kg}$. Đúng.
    \item Lần 2 nhiều hơn lần 1. Ý 3 Sai, Ý 4 Đúng.
\end{itemize}
\begin{itemize}
    \item a) \textbf{Đúng}.
    \item b) \textbf{Đúng}.
    \item c) \textbf{Sai}.
    \item d) \textbf{Đúng}.
\end{itemize}
}
\end{ex}

\textbf{Phần 3. Trả lời ngắn}
\setcounterpageref{ex}{0}

\begin{ex} %[3D16H81]
Tính: $42\,000 + 38\,000$.
\shortans{80000}
\loigiai{
\textbf{Đáp số:} $80000$\\
$42000 + 38000 = 80000$.
}
\end{ex}

\begin{ex} %[3D16H81]
Tính: $100\,000 - 36\,000$.
\shortans{64000}
\loigiai{
\textbf{Đáp số:} $64000$\\
$100000 - 36000 = 64000$.
}
\end{ex}

\begin{ex} %[3D16V81]
Tính: $16\,000 \times 5$.
\shortans{80000}
\loigiai{
\textbf{Đáp số:} $80000$\\
$16000 \times 5 = 80000$.
}
\end{ex}

\begin{ex} %[3D16V81]
Tính: $96\,000 : 3$.
\shortans{32000}
\loigiai{
\textbf{Đáp số:} $32000$\\
$96000 : 3 = 32000$.
}
\end{ex}

\begin{ex} %[3H16V81]
Tính diện tích hình chữ nhật có chiều dài $18\text{ cm}$, chiều rộng $5\text{ cm}$ theo xăng-ti-mét vuông.
\shortans{90}
\loigiai{
\textbf{Đáp số:} $90$\\
$18 \times 5 = 90\text{ cm}^2$.
}
\end{ex}

\begin{ex} %[3D16C81]
Tính giá trị biểu thức: $15\,000 \times 4 + 80\,000 : 2$.
\shortans{100000}
\loigiai{
\textbf{Đáp số:} $100000$\\
$60000 + 40000 = 100\,000$.
}
\end{ex}
